% ---- score2 standard preamble (from: python3 score2_tools.py preamble) ----
\documentclass[10pt]{extarticle}
\usepackage[a4paper,margin=1.5cm]{geometry}
\usepackage[T1]{fontenc}
\usepackage{amsmath,amssymb}
\usepackage{graphicx}
\usepackage[export]{adjustbox}
\usepackage{multicol}
\usepackage{enumitem}
\usepackage{booktabs,array,tabularx}
\setlength{\columnsep}{0.7cm}
\setlength{\parindent}{0pt}
\setlength{\parskip}{3pt}
\pagestyle{empty}
\makeatletter
% figure registry (figdefs.tex is generated by the extract step)
\newcommand{\qfigdef}[4]{\@namedef{qf@file@#1}{#2}\@namedef{qf@w@#1}{#3}\@namedef{qf@wide@#1}{#4}}
\newcommand{\qf@check}[1]{\@ifundefined{qf@file@#1}{\PackageError{score2}{Unknown figure id #1}{}}{}}
\newcommand{\qf@img}[4][]{\includegraphics[width=#2\dimexpr\@nameuse{qf@w@#4}pt\relax,max width=#3,max height=0.45\textheight,#1]{\@nameuse{qf@file@#4}}}
% Full-width questions: finalize wraps every question that contains a wide figure in \qfullbegin...\qfullend,
% so the WHOLE question (stem, figures, options) is set across the page in reading order, then the two columns
% resume. The question's head - its number, stem and first figure (\qfullsplit marks the end of the head) -
% is kept on one page, so a number is never left at a page foot with its figure overleaf; the rest may break
% across pages normally, so no big gaps are left. Without a \qfullsplit, the whole question is kept together
% if it is under 45% of a page. Never write these yourself.
\newif\ifqf@full
\newif\ifqf@boxed
\newsavebox{\qf@box}
\newcommand{\qf@open}{\global\qf@fulltrue\par\global\qf@boxedtrue
  \setbox\qf@box\vbox\bgroup\hsize\textwidth\linewidth\textwidth}
\newcommand{\qf@emit}[1]{\par\egroup\global\qf@boxedfalse
  \ifdim\dimexpr\ht\qf@box+\dp\qf@box\relax<#1\textheight\noindent\usebox{\qf@box}\par\else\unvbox\qf@box\fi}
\newcommand{\qfullsplit}{\ifqf@boxed\qf@emit{0.85}\fi}
\newcommand{\qf@close}{\ifqf@boxed\qf@emit{0.45}\fi\global\qf@fullfalse}
\newcommand{\qfullbegin}{\end{multicols}\qf@open}
\newcommand{\qfullend}{\qf@close\begin{multicols}{2}}
\newcommand{\qfullstart}{\qf@open}
\newcommand{\qfulldone}{\qf@close}
% \qfig[scale]{id}: block figure, centred, directly where it is written (size set by the tool).
% Fallback if a wide figure is not inside a full-width question: break out of the columns in place - no floats.
\newcommand{\qfig}[2][1]{\qf@check{#2}%
  \ifqf@full\qf@blk{#1}{#2}\else
  \ifnum\@nameuse{qf@wide@#2}=1 \expandafter\@firstoftwo\else\expandafter\@secondoftwo\fi
  {\end{multicols}\begin{center}\qf@img{#1}{\textwidth}{#2}\end{center}\begin{multicols}{2}}%
  {\qf@blk{#1}{#2}}\fi}
\newcommand{\qf@blk}[2]{\par\smallskip{\centering\qf@img{#1}{\linewidth}{#2}\par}\smallskip}
% \qfigi[scale]{id}: inline figure (inside option lists / tables / minipages); never breaks columns.
\newcommand{\qfigi}[2][1]{\qf@check{#2}\qf@img[valign=c]{#1}{\linewidth}{#2}}
\makeatother
% structure macros - lint reads these to check completeness and figure placement
\newcommand{\papersec}[1]{\par\bigskip{\centering\Large\bfseries PAPER-#1\par}\medskip}
\newcommand{\partsec}[2]{\par\medskip\noindent{\large\bfseries PART-#1: #2}\par\smallskip}
\newcommand{\sect}[1]{\par\smallskip\noindent\textbf{#1}\par}
\newcommand{\qn}[1]{\par\smallskip\noindent\textbf{#1.}~}
\newcommand{\paraq}[2]{\par\smallskip\noindent\textbf{Paragraph for Questions #1 and #2}\par}
\newenvironment{qoptions}{\begin{itemize}[leftmargin=1.8em,itemsep=1pt,topsep=2pt,label={}]}{\end{itemize}}
\newcommand{\opt}[2]{\item[(#1)] #2}
% figure registry (generated by score2_tools.py finalize)
\qfigdef{p1_phy_s1_q01_1}{TOAS - Test 03 (27 Feb 2022)-img/p1_phy_s1_q01_1.png}{124.6}{0}
\qfigdef{p1_phy_s1_q03_1}{TOAS - Test 03 (27 Feb 2022)-img/p1_phy_s1_q03_1.png}{108.7}{0}
\qfigdef{p1_phy_s1_q04_1}{TOAS - Test 03 (27 Feb 2022)-img/p1_phy_s1_q04_1.png}{45.6}{0}
\qfigdef{p1_phy_s1_q05_1}{TOAS - Test 03 (27 Feb 2022)-img/p1_phy_s1_q05_1.png}{122.5}{0}
\qfigdef{p1_phy_s1_q06_1}{TOAS - Test 03 (27 Feb 2022)-img/p1_phy_s1_q06_1.png}{124.9}{0}
\qfigdef{p1_phy_s2_q08_1}{TOAS - Test 03 (27 Feb 2022)-img/p1_phy_s2_q08_1.png}{210.6}{0}
\qfigdef{p1_phy_s2_q09_1}{TOAS - Test 03 (27 Feb 2022)-img/p1_phy_s2_q09_1.png}{113.4}{0}
\qfigdef{p1_phy_s2_q11_1}{TOAS - Test 03 (27 Feb 2022)-img/p1_phy_s2_q11_1.png}{163.1}{0}
\qfigdef{p1_phy_s3_q01_1}{TOAS - Test 03 (27 Feb 2022)-img/p1_phy_s3_q01_1.png}{99.6}{0}
\qfigdef{p1_phy_s3_q02_1}{TOAS - Test 03 (27 Feb 2022)-img/p1_phy_s3_q02_1.png}{66.4}{0}
\qfigdef{p1_phy_s3_q03_1}{TOAS - Test 03 (27 Feb 2022)-img/p1_phy_s3_q03_1.png}{128.6}{0}
\qfigdef{p1_phy_s3_q05_1}{TOAS - Test 03 (27 Feb 2022)-img/p1_phy_s3_q05_1.png}{44.3}{0}
\qfigdef{p1_chem_s1_q01_A}{TOAS - Test 03 (27 Feb 2022)-img/p1_chem_s1_q01_A.png}{41.3}{0}
\qfigdef{p1_chem_s1_q01_B}{TOAS - Test 03 (27 Feb 2022)-img/p1_chem_s1_q01_B.png}{81.1}{0}
\qfigdef{p1_chem_s1_q01_C}{TOAS - Test 03 (27 Feb 2022)-img/p1_chem_s1_q01_C.png}{53.2}{0}
\qfigdef{p1_chem_s1_q01_D}{TOAS - Test 03 (27 Feb 2022)-img/p1_chem_s1_q01_D.png}{123.0}{0}
\qfigdef{p1_chem_s1_q02_A}{TOAS - Test 03 (27 Feb 2022)-img/p1_chem_s1_q02_A.png}{80.6}{0}
\qfigdef{p1_chem_s1_q02_B}{TOAS - Test 03 (27 Feb 2022)-img/p1_chem_s1_q02_B.png}{94.8}{0}
\qfigdef{p1_chem_s1_q02_C}{TOAS - Test 03 (27 Feb 2022)-img/p1_chem_s1_q02_C.png}{161.1}{0}
\qfigdef{p1_chem_s1_q02_D}{TOAS - Test 03 (27 Feb 2022)-img/p1_chem_s1_q02_D.png}{65.9}{0}
\qfigdef{p1_chem_s1_q04_1}{TOAS - Test 03 (27 Feb 2022)-img/p1_chem_s1_q04_1.png}{119.0}{0}
\qfigdef{p1_chem_s2_q07_A}{TOAS - Test 03 (27 Feb 2022)-img/p1_chem_s2_q07_A.png}{110.8}{0}
\qfigdef{p1_chem_s2_q07_B}{TOAS - Test 03 (27 Feb 2022)-img/p1_chem_s2_q07_B.png}{173.6}{0}
\qfigdef{p1_chem_s2_q07_C}{TOAS - Test 03 (27 Feb 2022)-img/p1_chem_s2_q07_C.png}{179.2}{0}
\qfigdef{p1_chem_s2_q07_D}{TOAS - Test 03 (27 Feb 2022)-img/p1_chem_s2_q07_D.png}{180.0}{0}
\qfigdef{p1_chem_s2_q08_1}{TOAS - Test 03 (27 Feb 2022)-img/p1_chem_s2_q08_1.png}{125.6}{0}
\qfigdef{p1_chem_s2_q08_2}{TOAS - Test 03 (27 Feb 2022)-img/p1_chem_s2_q08_2.png}{46.3}{0}
\qfigdef{p1_chem_s2_q08_3}{TOAS - Test 03 (27 Feb 2022)-img/p1_chem_s2_q08_3.png}{46.3}{0}
\qfigdef{p1_chem_s2_q09_1}{TOAS - Test 03 (27 Feb 2022)-img/p1_chem_s2_q09_1.png}{90.8}{0}
\qfigdef{p1_chem_s3_q01_1}{TOAS - Test 03 (27 Feb 2022)-img/p1_chem_s3_q01_1.png}{105.3}{0}
\qfigdef{p2_phy_s1_q01_1}{TOAS - Test 03 (27 Feb 2022)-img/p2_phy_s1_q01_1.png}{92.5}{0}
\qfigdef{p2_phy_s1_q02_1}{TOAS - Test 03 (27 Feb 2022)-img/p2_phy_s1_q02_1.png}{177.5}{0}
\qfigdef{p2_phy_s1_q03_1}{TOAS - Test 03 (27 Feb 2022)-img/p2_phy_s1_q03_1.png}{119.5}{0}
\qfigdef{p2_phy_s1_q04_1}{TOAS - Test 03 (27 Feb 2022)-img/p2_phy_s1_q04_1.png}{152.5}{0}
\qfigdef{p2_phy_s1_q05_1}{TOAS - Test 03 (27 Feb 2022)-img/p2_phy_s1_q05_1.png}{182.5}{0}
\qfigdef{p2_phy_s1_q06_1}{TOAS - Test 03 (27 Feb 2022)-img/p2_phy_s1_q06_1.png}{140.6}{0}
\qfigdef{p2_phy_s2_q02_1}{TOAS - Test 03 (27 Feb 2022)-img/p2_phy_s2_q02_1.png}{166.0}{0}
\qfigdef{p2_phy_s2_q03_1}{TOAS - Test 03 (27 Feb 2022)-img/p2_phy_s2_q03_1.png}{124.2}{0}
\qfigdef{p2_phy_s2_q04_1}{TOAS - Test 03 (27 Feb 2022)-img/p2_phy_s2_q04_1.png}{168.9}{0}
\qfigdef{p2_phy_s2_q05_1}{TOAS - Test 03 (27 Feb 2022)-img/p2_phy_s2_q05_1.png}{174.2}{0}
\qfigdef{p2_phy_s2_q06_1}{TOAS - Test 03 (27 Feb 2022)-img/p2_phy_s2_q06_1.png}{201.5}{0}
\qfigdef{p2_phy_s3_q01_1}{TOAS - Test 03 (27 Feb 2022)-img/p2_phy_s3_q01_1.png}{166.8}{0}
\qfigdef{p2_chem_s1_q01_A}{TOAS - Test 03 (27 Feb 2022)-img/p2_chem_s1_q01_A.png}{43.0}{0}
\qfigdef{p2_chem_s1_q01_B}{TOAS - Test 03 (27 Feb 2022)-img/p2_chem_s1_q01_B.png}{38.4}{0}
\qfigdef{p2_chem_s1_q01_C}{TOAS - Test 03 (27 Feb 2022)-img/p2_chem_s1_q01_C.png}{38.2}{0}
\qfigdef{p2_chem_s1_q01_D}{TOAS - Test 03 (27 Feb 2022)-img/p2_chem_s1_q01_D.png}{38.7}{0}
\qfigdef{p2_chem_s1_q02_A}{TOAS - Test 03 (27 Feb 2022)-img/p2_chem_s1_q02_A.png}{175.6}{0}
\qfigdef{p2_chem_s1_q02_2}{TOAS - Test 03 (27 Feb 2022)-img/p2_chem_s1_q02_2.png}{142.1}{0}
\qfigdef{p2_chem_s1_q02_3}{TOAS - Test 03 (27 Feb 2022)-img/p2_chem_s1_q02_3.png}{106.5}{0}
\qfigdef{p2_chem_s1_q02_D}{TOAS - Test 03 (27 Feb 2022)-img/p2_chem_s1_q02_D.png}{183.2}{0}
\qfigdef{p2_chem_s1_q06_1}{TOAS - Test 03 (27 Feb 2022)-img/p2_chem_s1_q06_1.png}{265.2}{1}
\qfigdef{p2_chem_s2_q01_1}{TOAS - Test 03 (27 Feb 2022)-img/p2_chem_s2_q01_1.png}{239.9}{0}
\qfigdef{p2_chem_s2_q02_1}{TOAS - Test 03 (27 Feb 2022)-img/p2_chem_s2_q02_1.png}{198.5}{0}
\qfigdef{p2_chem_s2_q06_1}{TOAS - Test 03 (27 Feb 2022)-img/p2_chem_s2_q06_1.png}{134.5}{0}
\qfigdef{p2_chem_s2_q06_2}{TOAS - Test 03 (27 Feb 2022)-img/p2_chem_s2_q06_2.png}{134.2}{0}
\qfigdef{p2_chem_s2_q06_3}{TOAS - Test 03 (27 Feb 2022)-img/p2_chem_s2_q06_3.png}{134.4}{0}
\qfigdef{p2_chem_s3_q02_1}{TOAS - Test 03 (27 Feb 2022)-img/p2_chem_s3_q02_1.png}{135.4}{0}
\raggedcolumns
\makeatletter\@beginparpenalty=10000\makeatother
\newcolumntype{L}[1]{>{\raggedright\arraybackslash}p{#1}}
\begin{document}
\begin{center}{\Large\bfseries ALLEN TOAS --- Test 03 (27 Feb 2022)}\end{center}
\begin{multicols}{2}
\papersec{1}
\partsec{1}{PHYSICS}
\sect{SECTION-I(i)}
\qn{1} Find the lateral magnification produced by the combination of lenses shown in the figure.
\qfig{p1_phy_s1_q01_1}
\begin{qoptions}
\opt{A}{2}
\opt{B}{4}
\opt{C}{6}
\opt{D}{8}
\end{qoptions}
\qn{2} A body is executing simple harmonic motion. At a displacement $x$, its potential energy is $E_1$ and at a displacement of $y$, its potential energy is $E_2$. The potential energy $E$ at a displacement $(x+y)$ is :-
\begin{qoptions}
\opt{A}{$E_1+E_2$}
\opt{B}{$\sqrt{E_1^2+E_2^2}$}
\opt{C}{$E_1+E_2+2\sqrt{E_1E_2}$}
\opt{D}{$\sqrt{E_1E_2}$}
\end{qoptions}
\qn{3} The minimum value of $\theta$ so that the ray gets totally reflected from spherical surface of radius of curvature $R$ is $\left[\dfrac{\mu_2}{\mu_1}=\dfrac12\right]$
\qfig{p1_phy_s1_q03_1}
\begin{qoptions}
\opt{A}{$30^\circ$}
\opt{B}{$60^\circ$}
\opt{C}{zero}
\opt{D}{not possible}
\end{qoptions}
\qn{4} In the figure shown, blocks $P$ and $Q$ are in contact but do not stick to each other. The springs are in their natural lengths. The system is released from rest. Then the distance between $P$ and $Q$ when $Q$ is at the lowest point for first time will be :-
\qfig{p1_phy_s1_q04_1}
\begin{qoptions}
\opt{A}{$\dfrac{2mg}{K}$}
\opt{B}{$\dfrac{4mg}{K}$}
\opt{C}{$\dfrac{3mg}{K}$}
\opt{D}{0}
\end{qoptions}
\qn{5} A block is placed on an inclined plane moving towards right horizontally with an acceleration $a_0=g$. The length of the plane $AC=1$ m. Friction is absent everywhere. The time taken by the block to reach from $C$ to $A$ is ($g=10\ \mathrm{m/s^2}$) :-
\qfig{p1_phy_s1_q05_1}
\begin{qoptions}
\opt{A}{1.2 s}
\opt{B}{0.74 s}
\opt{C}{2.56 s}
\opt{D}{0.42 s}
\end{qoptions}
\qn{6} The graph shows the variation of $\dfrac1V$ (where $V$ is the velocity of the particle) with respect to time. Then find the value of acceleration at $t=3$ sec.
\qfig{p1_phy_s1_q06_1}
\begin{qoptions}
\opt{A}{$5\ \mathrm{m/s^2}$}
\opt{B}{$1\ \mathrm{m/s^2}$}
\opt{C}{$3\sqrt3\ \mathrm{m/s^2}$}
\opt{D}{$3\ \mathrm{m/s^2}$}
\end{qoptions}
\sect{SECTION-I(ii)}
\qn{7} Consider the linear graph between magnification ($m$) and image position ($v$) for a lens. Which of the following statement about the straight line is correct ?
\begin{qoptions}
\opt{A}{Focal length of the lens is equal to intercept on $v$-axis}
\opt{B}{Magnitude of focal length of the lens is equal to inverse of slope of the line.}
\opt{C}{Magnitude of intercept on $m$-axis is equal to unity}
\opt{D}{None of these}
\end{qoptions}
\qn{8} Kinetic energy of a particle executing SHM changes with time according to following graph. Mass of the particle is 2 kg. Then :-
\qfig{p1_phy_s2_q08_1}
\begin{qoptions}
\opt{A}{amplitude of oscillation is 0.4 m}
\opt{B}{acceleration of the particle is zero at $t=0.1$ s}
\opt{C}{time period of oscillation is 0.1 s}
\opt{D}{frequency of oscillation is 5 Hz.}
\end{qoptions}
\qn{9} An object is moving in a circular path of radius 1 m around (0, 10 m) in $xy$ plane with constant speed 1 m/s. Speed of image w.r.t. object is $v$.
\qfig{p1_phy_s2_q09_1}
\begin{qoptions}
\opt{A}{$v_{\max}=2$ m/s}
\opt{B}{$v_{\min}=1$ m/s}
\opt{C}{$v$ is minimum when object is closest or farthest from the mirror}
\opt{D}{$v$ is maximum when object crosses $y$--axis}
\end{qoptions}
\qn{10} A particle of mass $m=3$ kg is moving on a straight line in a conservative force field in which the potential energy associated with the particle is given by $U(y)=2y^2-30y+450$. At $t=0$, particle is at $y=8$ m and moving towards positive $y$-axis with 1 m/s.
\begin{qoptions}
\opt{A}{During the course of its motion the maximum speed of the particle will be $\sqrt{\dfrac43}$ m/s}
\opt{B}{Particle will never cross the origin.}
\opt{C}{Particle will come to its equilibrium position for the first time at $\dfrac{5\pi}{4\sqrt3}$ second.}
\opt{D}{The amplitude of oscillation for the particle is 1 m.}
\end{qoptions}
\qn{11} A wedge is moving rightwards on which a block of mass 10 kg is placed on it. Friction coefficient between the wedge and the block is 0.8. [take $g=10\ \mathrm{m/s^2}$]. Select correct alternative(s) among the following options.
\qfig{p1_phy_s2_q11_1}
\begin{qoptions}
\opt{A}{If wedge is moving with constant velocity then acceleration of block is zero.}
\opt{B}{If wedge is moving with $\vec a=2\left(\hat i\right)\ \mathrm{m/s^2}$ then friction acting on block is 44N.}
\opt{C}{If wedge is moving with constant velocity then friction acting on block is 64N.}
\opt{D}{If wedge is moving with $\vec a=10\left(\hat i\right)\ \mathrm{m/s^2}$ then friction is 20N, downward on the block along the incline.}
\end{qoptions}
\qn{12} The instantaneous velocity $v$ of a particle is related to its displacement $x$ according to the relation $v=8x+1$ (All parameters in SI units). Which of the following statement(s) is (are) true if $x=0$ at $t=0$ :-
\begin{qoptions}
\opt{A}{The particle will experience a retardation.}
\opt{B}{The particle will be at rest at $t=0$.}
\opt{C}{The displacement of the particle at time $t$ is $x=\dfrac18\left(e^{8t}-1\right)$}
\opt{D}{The particle will experience an acceleration.}
\end{qoptions}
\sect{SECTION-II}
\qn{1} A man `A' stands at the position shown in the figure and a second man `B' approaches the mirror along the line perpendicular to it which passes through its centre. At what distance (in meter) from the mirror will `B' be at the moment when `A' and `B' first see each other in the mirror ?
\qfig{p1_phy_s3_q01_1}
\qn{2} A uniform solid hemisphere of mass $m$ and radius $R$ is attached to the roof with a chord of torsional constant $C$ and performing torsional SHM. Then the time period (in seconds) of SHM is (Take $m=15$ kg, $R=\dfrac2\pi$ m, $C=6$ Nm/rad.)
\qfig{p1_phy_s3_q02_1}
\qn{3} A ray is incident on a glass sphere as shown. The opposite surface of the sphere is partially silvered. If the net deviation of the ray transmitted at the partially silvered surface is $1/3$rd of the net deviation suffered by the ray coming out after one reflection, the refractive index of the sphere is $n$. Find the value of $n$.
\qfig{p1_phy_s3_q03_1}
\qn{4} Two particles are performing SHM with same amplitude and time period. At an instant two particles are having velocity 1 m/s but one is on the right and the other is on the left of their mean position. When the particles have same position, their speed is $\sqrt3$ m/s. Find the maximum speed (in m/s) of particles during SHM.
\qn{5} In the figure shown, pulley and spring are ideal and strings are light and inextensible. Initially all the bodies are at rest when string connecting $A$ and $B$ is cut. Find initial acceleration of $C$.
\qfig{p1_phy_s3_q05_1}
\qn{6} Two particles $A$ and $B$ located at points $(0,-10\sqrt3)$ and $(0,0)$ respectively in $xy$ plane. They start moving simultaneously at time $t=0$ with constant velocities $\vec V_A=5\hat i$ m/s and $\vec V_B=-5\sqrt3\,\hat j$ m/s, respectively. Time when they are closest to each other is found to be $K$ second. Find $K$. All distance are given in meter.
\partsec{2}{CHEMISTRY}
\sect{SECTION-I(i)}
\qn{1} Which of the following molecule is resolvable / optically active?
\begin{qoptions}
\opt{A}{\qfigi{p1_chem_s1_q01_A}}
\opt{B}{\qfigi{p1_chem_s1_q01_B}}
\opt{C}{\qfigi{p1_chem_s1_q01_C}}
\opt{D}{\qfigi{p1_chem_s1_q01_D}}
\end{qoptions}
\qn{2} Which of following gives $2^\circ$ alcohol major product?
\begin{qoptions}
\opt{A}{\qfigi{p1_chem_s1_q02_A}}
\opt{B}{\qfigi{p1_chem_s1_q02_B}}
\opt{C}{\qfigi{p1_chem_s1_q02_C}}
\opt{D}{\qfigi{p1_chem_s1_q02_D}}
\end{qoptions}
\qn{3} Assuming 100\% ionization, the increasing order of the freezing point of the solution will be :-
\begin{qoptions}
\opt{A}{$0.10\ \mathrm{mol\,kg^{-1}\ Ba_3(PO_4)_2}$ $<$ $0.10\ \mathrm{mol\,kg^{-1}\ Na_2SO_4}$ $<$ $0.10\ \mathrm{mol\,kg^{-1}\ KCl}$}
\opt{B}{$0.10\ \mathrm{mol\,kg^{-1}\ KCl}$ $<$ $0.10\ \mathrm{mol\,kg^{-1}\ Na_2SO_4}$ $<$ $0.10\ \mathrm{mol\,kg^{-1}\ Ba_3(PO_4)_2}$}
\opt{C}{$0.10\ \mathrm{mol\,kg^{-1}\ Na_2SO_4}$ $<$ $0.10\ \mathrm{mol\,kg^{-1}\ Ba_3(PO_4)_2}$ $<$ $0.10\ \mathrm{mol\,kg^{-1}\ KCl}$}
\opt{D}{$0.10\ \mathrm{mol\,kg^{-1}\ KCl}$ $<$ $0.10\ \mathrm{mol\,kg^{-1}\ Ba_3(PO_4)_2}$ $<$ $0.10\ \mathrm{mol\,kg^{-1}\ Na_2SO_4}$}
\end{qoptions}
\qn{4} Correct expression for determining packing fraction (P.F.) of an ideal NaCl unit cell, if all ions along an edge diagonal are absent, as shown [$\mathrm{Cl^-}$ is forming FCC, $r_+$ = Cationic radius, $r_-$ = Anionic radius]
\qfig{p1_chem_s1_q04_1}
\begin{qoptions}
\opt{A}{$\dfrac{\frac43\pi\left(r_+^3+r_-^3\right)}{16\sqrt2\,r_-^3}$}
\opt{B}{$\dfrac{\frac43\pi\left(\frac72\cdot r_+^3+4r_-^3\right)}{16\sqrt2\,r_-^3}$}
\opt{C}{$\dfrac{\frac43\pi\left(\frac52 r_+^3+4r_-^3\right)}{16\sqrt2\,r_-^3}$}
\opt{D}{$\dfrac{\frac43\pi\left(\frac52 r_+^3+4r_-^3\right)}{16\sqrt2\,r_+^3}$}
\end{qoptions}
\qn{5} Which of the following statement is correct?
\begin{qoptions}
\opt{A}{for $\mathrm{N_2\to N_2^+}$ bond order decreases and magnetic nature does not change.}
\opt{B}{for $\mathrm{CO\to CO^+}$ bond order decreases and magnetic nature changes.}
\opt{C}{for $\mathrm{O_2\to O_2^{2-}}$ spin multiplicity decreases and magnetic nature changes.}
\opt{D}{for $\mathrm{B_2\to B_2^+}$ bond order increases and magnetic nature does not changes.}
\end{qoptions}
\qn{6} In which of the following maximum difference of \% s-character between hybrid orbital of cationic and anionic part.
\begin{qoptions}
\opt{A}{$\mathrm{Cl_2O_{6(s)}}$}
\opt{B}{$\mathrm{N_2O_{5(s)}}$}
\opt{C}{$\mathrm{I_2Cl_{6(\ell)}}$}
\opt{D}{$\mathrm{PCl_{5(s)}}$}
\end{qoptions}
\sect{SECTION-I(ii)}
\qn{7} Which of following option is/are correct?
\begin{qoptions}
\opt{A}{\qfigi{p1_chem_s2_q07_A}}
\opt{B}{\qfigi{p1_chem_s2_q07_B}}
\opt{C}{\qfigi{p1_chem_s2_q07_C}}
\opt{D}{\qfigi{p1_chem_s2_q07_D}}
\end{qoptions}
\qn{8} \qfigi{p1_chem_s2_q08_1}

Correct option(s) for above reaction sequence
\begin{qoptions}
\opt{A}{(P) is \qfigi{p1_chem_s2_q08_2}}
\opt{B}{(Q) is \qfigi{p1_chem_s2_q08_3}}
\opt{C}{On increasing the temperature \% of (Q) increases in product mixture}
\opt{D}{Product (P) is formed by more stable carbocation formation}
\end{qoptions}
\qn{9} Study the following figure, and select the correct statements (assuming complete dissociation of the electrolyte).
\qfig{p1_chem_s2_q09_1}
\begin{qoptions}
\opt{A}{There will be no movement of any solution across the membrane.}
\opt{B}{Water from $\mathrm{BaCl_2}$ will flow towards the NaCl solution.}
\opt{C}{Water from NaCl will flow towards the $\mathrm{BaCl_2}$.}
\opt{D}{The osmotic pressure of 0.1 M NaCl is higher than the osmotic pressure of 0.05 M $\mathrm{BaCl_2}$,}
\end{qoptions}
\qn{10} Select the correct statement(s) :
\begin{qoptions}
\opt{A}{The ionic crystal of AgBr may have Schottky defect}
\opt{B}{The unit cell having crystal parameters, $a=b\neq c$, $\alpha=\beta=90^\circ$, $\gamma=120^\circ$ is hexagonal}
\opt{C}{In ionic compounds having Frenkel defect, the ratio $r^+/r^-$ is high}
\opt{D}{The co-ordination number of $\mathrm{Na^+}$ ion in NaCl is 6}
\end{qoptions}
\qn{11} Between $\mathrm{ClO_2}$ and $\mathrm{Cl_2O}$
\begin{qoptions}
\opt{A}{The Cl - O bond length would be larger in $\mathrm{Cl_2O}$}
\opt{B}{The bond angle would be larger in $\mathrm{Cl_2O}$}
\opt{C}{$\mathrm{ClO_2}$ is likely to be a coloured compound}
\opt{D}{$\mathrm{Cl_2O}$ is likely to be soluble in water}
\end{qoptions}
\qn{12} Which of the following can be explained on the basis of Fajan's Rules.
\begin{qoptions}
\opt{A}{$\mathrm{Ag_2S}$ is less soluble than $\mathrm{Ag_2O}$.}
\opt{B}{$\mathrm{Fe(OH)_3}$ is less soluble than $\mathrm{Fe(OH)_2}$.}
\opt{C}{$\mathrm{BaCO_3}$ is less soluble than $\mathrm{MgCO_3}$.}
\opt{D}{$\mathrm{BeCO_3}$ is thermally less stable than $\mathrm{BaCO_3}$}
\end{qoptions}
\sect{SECTION-II}
\qn{1} How many chiral center are present in this compound
\qfig{p1_chem_s3_q01_1}
\qn{2} An unknown unbranched hydrocarbon with molecular formula $\mathrm{C_5H_{10}}$ decolourise $\mathrm{Br_2+H_2O}$ and on reaction with H--Cl/$\mathrm{CCl_4}$ give 2-chloro pentane as major product. Identify total number of possible reactants (excluding stereochemistry) -
\qn{3} A 1 litre solution containing equal moles of FeO and $\mathrm{Fe_{0.8}O}$ was titrated with 70 ml 0.3M $\mathrm{KMnO_4}$ in acidic medium. Millimoles of $\mathrm{Fe^{3+}}$ produced are -
\qn{4} If shortest distance between two nearest $\mathrm{Zn^{2+}}$ ion in ZnS crystal is $3\sqrt2$ \AA. Then edge length of this crystal (in \AA) will be
\qn{5} Find the no. of species having fractional bond order and paramagnetic in nature?

$\mathrm{N_2^+,\ O_2,\ O_2^+,\ F_2,\ B_2^+,\ C_2^+,\ CN^-}$
\qn{6} $\mathrm{XY_n}$ is an interhalogen compound which produces oxyacid of $+5$ oxidation state of central atom on complete hydrolysis. If `P' moles $\mathrm{H_2O}$ is required for the complete hydrolysis of one mole $\mathrm{XY_n}$ then the value of $(n+P)$ is :
\partsec{3}{MATHEMATICS}
\sect{SECTION-I(i)}
\qn{1} Number of possible ordered pair(s) $(x,y)$ satisfying the equation $\cot^{-1}x=-e^{y}$ is -
\begin{qoptions}
\opt{A}{0}
\opt{B}{1}
\opt{C}{2}
\opt{D}{3}
\end{qoptions}
\qn{2} The complete set of values of $x$ for which $|\tan^{-1}x|+|\cot^{-1}x|=\dfrac\pi2$ is valid
\begin{qoptions}
\opt{A}{$(-3,3)$}
\opt{B}{$(-1,1)$}
\opt{C}{$(-\infty,0]$}
\opt{D}{$[0,\infty)$}
\end{qoptions}
\qn{3} Let $f:\mathbb R\to\mathbb R$ be a non-constant continuous function such that $f(2x)=(e^x+1)f(x)$, $f(0)=0$ \& $f'(0)=1$, then range of $f(x)$ is-
\begin{qoptions}
\opt{A}{$(-1,\infty)$}
\opt{B}{$[0,\infty)$}
\opt{C}{$[1,\infty)$}
\opt{D}{$(1,\infty)$}
\end{qoptions}
\qn{4} The greatest value of the function $f(x)=(\sin x-\cos x)^2+(\sin x+1)^2+(\cos x-1)^2$ is :-
\begin{qoptions}
\opt{A}{$5+2\sqrt2$}
\opt{B}{$5-2\sqrt2$}
\opt{C}{24}
\opt{D}{16}
\end{qoptions}
\qn{5} If $\dfrac1{1^5}+\dfrac1{2^5}+\dfrac1{3^5}+\dfrac1{4^5}+\ldots$ upto infinite terms $=k$, then value of $\left(\dfrac1{1^5}+\dfrac1{3^5}+\dfrac1{5^5}+\dfrac1{7^5}+\ldots\right.$ upto infinite terms$\left.\right)$ will be-
\begin{qoptions}
\opt{A}{$\dfrac k{32}$}
\opt{B}{$\dfrac{31k}{32}$}
\opt{C}{$\dfrac k{16}$}
\opt{D}{$\dfrac{15k}{16}$}
\end{qoptions}
\qn{6} If sum of first $n$ terms of a series is $S_n=3n^3+2n^2$, then $6^{\mathrm{th}}$ term is-
\begin{qoptions}
\opt{A}{309}
\opt{B}{720}
\opt{C}{101}
\opt{D}{295}
\end{qoptions}
\sect{SECTION-I(ii)}
\qn{7} Let $S_n=(2+2)+(2^2+5)+(2^3+10)+(2^4+17)+\ldots$ upto $n$ brackets. If $S_n=2^{n+A}+Bn^3+Cn^2+Dn+E\ \forall\ n\in\mathbb N$ (where $A,B,C,D,E$ are constants), then which of the following is/are true ?
\begin{qoptions}
\opt{A}{$A=1$}
\opt{B}{$B=\dfrac16$}
\opt{C}{$D=\dfrac76$}
\opt{D}{$E=-2$}
\end{qoptions}
\qn{8} If $x$ is a non-zero real, then value of $x^2+\dfrac{27}{x^2}+1$ can be-
\begin{qoptions}
\opt{A}{8}
\opt{B}{9}
\opt{C}{13}
\opt{D}{20}
\end{qoptions}
\qn{9} If $f(x)=\sqrt{|x|^2-5|x|+6}+\sqrt{8+2|x|-|x|^2}$, then $x$ can be-
\begin{qoptions}
\opt{A}{$[-4,-3]$}
\opt{B}{$[-3,-2]$}
\opt{C}{$[-2,2]$}
\opt{D}{$[3,4]$}
\end{qoptions}
\qn{10} If $6\sin^{-1}\left(x^2-4x+\dfrac72\right)=\pi$, then-
\begin{qoptions}
\opt{A}{$x=1$}
\opt{B}{$x=2$}
\opt{C}{$x=3$}
\opt{D}{$x=4$}
\end{qoptions}
\qn{11} The algebraic expression for \[f(x)=\tan\left(\sin^{-1}\left(\cos\left(\tan^{-1}\dfrac x2\right)\right)\right)\] can be-
\begin{qoptions}
\opt{A}{$\dfrac\pi2$}
\opt{B}{$\dfrac1x$}
\opt{C}{$\dfrac2x$}
\opt{D}{$-\dfrac2x$}
\end{qoptions}
\qn{12} Which of the following is (are) positive -
\begin{qoptions}
\opt{A}{$\log_7 25$}
\opt{B}{$\log_3(\log_3 9)$}
\opt{C}{$\dfrac{\log_3 5}{\log_3(2+x^2)}$}
\opt{D}{$\left(\log_{1/2}3\right)^{10}$}
\end{qoptions}
\sect{SECTION-II}
\qn{1} Let $f$ be a function satisfying the functional equation $f(x)+2f\left(\dfrac{2x+1}{x-2}\right)=3x,\ x\neq2$, then the value of $\dfrac12\left[\left|\dfrac{f(3)}{2f(7)}\right|\right]$ is [Note : [.] denotes greatest integer function]
\qn{2} The number of integers in the domain of $f(x)=\dfrac{\sqrt{(x^2+kx+k)}}{x^2+k}$ is
\qn{3} If $x$ and $y$ are positive integers satisfying \[\tan^{-1}\left(\dfrac1x\right)+\tan^{-1}\left(\dfrac1y\right)=\tan^{-1}\left(\dfrac17\right),\] then the number of ordered pairs of $(x,y)$ is
\qn{4} If $\theta=\sin^{-1}\left(\dfrac{2x}{1+x^2}\right)+\cos^{-1}\left(\dfrac{1-x^2}{1+x^2}\right)$ for $x\geqslant\dfrac32$, then the absolute value of $\dfrac12\left(\dfrac{\cos\theta+\tan\theta+4}{\sec\theta}\right)$ is
\qn{5} If $\sec(2\theta-2\alpha),\ \dfrac12\mathrm{cosec}^2\theta,\ \sec(2\theta+2\alpha)$ are in H.P., then the value of $\sec2\theta\sec^2\alpha$ is
\qn{6} Let $a,b,c$ are real number forming an A.P. and $(1+a),(2+b),(1+c)$ are in G.P., then minimum value of $ac$ is
\newpage
\papersec{2}
\partsec{1}{PHYSICS}
\sect{SECTION-I}
\qn{1} A small block of mass $m$ is fixed at upper end of a massless vertical spring of spring constant $K=\dfrac{4mg}{L}$ and natural length 10 m. The lower end of spring is free and is at a height $L$ meter from fixed horizontal floor as shown. The spring is initially unstretched and the spring block system is released from rest in the shown position.
\qfig{p2_phy_s1_q01_1}
\begin{qoptions}
\opt{A}{When the speed of block is maximum the magnitude of force exerted by spring on the block is $mg$.}
\opt{B}{The maximum speed attained by the block is $\dfrac32\sqrt{gL}$.}
\opt{C}{The maximum compression of the spring is $\dfrac L8$}
\opt{D}{The time taken by the mass to move from the equilibrium position to the maximum compressed position of the spring is $\dfrac\pi4\sqrt{\dfrac Lg}$}
\end{qoptions}
\qn{2} When object `O' moves towards a fixed lens mirror combination, select correct choice/choices:
\qfig{p2_phy_s1_q02_1}
\begin{qoptions}
\opt{A}{Image moves towards negative $x$-axis.}
\opt{B}{Speed of image and object may be same}
\opt{C}{Image may move faster than object}
\opt{D}{Image may come closer to arrangement}
\end{qoptions}
\qn{3} A block having mass $m$ is placed over a inclined plane of same mass and inclination $45^\circ$ if coefficient of friction between all surface is same it is $\mu$. Now system set free to move then
\qfig{p2_phy_s1_q03_1}
\begin{qoptions}
\opt{A}{For $\mu>1$ block over incline will move but wedge don't}
\opt{B}{For $0.2<\mu<1$ block will move but wedge will not}
\opt{C}{For $\mu=0.15$ both will move}
\opt{D}{It is not possible to move wedge without moving block over wedge}
\end{qoptions}
\qn{4} An equiconvex lens of focal length 20 cm having a glass of refractive index $\dfrac32$ is used in a medium air on left side and another medium having refractive index $\mu$ on the right side of lens as shown in figure.
\qfig{p2_phy_s1_q04_1}
\begin{qoptions}
\opt{A}{Character of lens will be diverging for $\mu>2$}
\opt{B}{Focal length of lens will be 40 cm for $\mu=4$ and character will be diverging}
\opt{C}{Object distance and image distance with respect to lens will be same for $\mu=2$}
\opt{D}{Character of lens can't be explained}
\end{qoptions}
\qn{5} $n$ number of identical equilateral prisms are kept in contact as shown in figure. If deviation through a single prism is $\delta$, then ($n$, $m$ are integers)
\qfig{p2_phy_s1_q05_1}
\begin{qoptions}
\opt{A}{if $n=2m$, deviation through $n$ prism is zero.}
\opt{B}{if $n=2m+1$, deviation through system of $n$ prisms is $\delta$.}
\opt{C}{if $n=2m$, deviation through system of $n$ prism is $\delta$.}
\opt{D}{if $n=2m+1$ deviation through system of $n$ prisms is zero}
\end{qoptions}
\qn{6} Consider an arrangement of two thin equiconvex lenses made of glass ($\mu_g=3/2$) and each of focal length in air 10 cm are separated by a distance 10 cm. The space between the two lenses is filled with water of refractive index $\mu_W=\dfrac43$. A small linear object `AB' of size 0.20 cm is placed perpendicular to the optic axis at a distance 10 cm from the first lens in air as shown. Then choose the correct option(s) :
\qfig{p2_phy_s1_q06_1}
\begin{qoptions}
\opt{A}{The position of the final image formed is 25 cm right to the second lens.}
\opt{B}{The position of the final image formed is 40 cm right to the second lens.}
\opt{C}{The size of the final image formed is 0.40 cm.}
\opt{D}{The size of the final image formed is 0.60 cm.}
\end{qoptions}
\sect{SECTION-II}
\qn{1} A projectile is thrown with a velocity of 18 m/s at an angle of $60^\circ$ with horizontal. The interval between the moments (in sec) when speed of the projectile is 15 m/s is ($g=10\ \mathrm{m/s^2}$)
\qn{2} The two ends $A$ and $B$ of a uniform rod of length $\ell=1$ m are moving with velocities $V_A$ and $V_B$ as shown. The velocity $V_B$ (m/s) is
\qfig{p2_phy_s2_q02_1}
\qn{3} A system of coordinates is drawn in a medium whose refractive index varies as $\mu=\dfrac2{1+y^2}$, where $0\leqslant y\leqslant1$. A ray of light is incident at origin at an angle $60^\circ$ with $y$-axis as shown in the figure. At point $P$ ray becomes parallel to $x$-axis. Find the value of `$H$'. (Take $\mu=2$ for $y<0$)
\qfig{p2_phy_s2_q03_1}
\qn{4} A commander fires a shell at certain angle of projection from `A' which clears the building (cuboid) of dimensions $20\times15\times12.5\ \mathrm{m^3}$ in its diagonal plane. Shell just clears the roof diagonally at points $P$ and $Q$ and falls on the other side of the building at $B$, and makes $45^\circ$ with horizontal at $P$ and $Q$ as shown in the figure. If the range $AB$ of the shell is $N\times\sqrt3$ metre then find $N$. ($g=10\ \mathrm{m/s^2}$)
\qfig{p2_phy_s2_q04_1}
\qn{5} In the figure shown below all the springs are ideal and time period of small oscillations about equilibrium position is $\pi\sqrt{\dfrac{xm}{k}}$. Find the value of $x$.
\qfig{p2_phy_s2_q05_1}
\qn{6} Four blocks are arranged on a smooth horizontal surface as shown. The masses of the blocks are given (see the figure). The coefficient of static friction between the top and the bottom blocks is $\mu_S$. What is the maximum value of the horizontal force $F$, applied to one of the bottom blocks as shown, that makes all the four blocks move with the same acceleration? [Take : $\mu_S=0.5$, $g=10\ \mathrm{m/s^2}$, $m=1$ kg, $M=2$ kg]
\qfig{p2_phy_s2_q06_1}
\sect{SECTION-III}
\qn{1} Distance between two images formed by upper and lower part of the point object placed at 30 cm from given lens is $(60+x)$ cm, then the value of $x$ is :
\qfig{p2_phy_s3_q01_1}
\qn{2} A block of mass 2 kg is gently placed over a massive plank moving horizontally over a smooth surface with velocity 6 m/s. The coefficient of friction between the block and plank is 0.2. The distance travelled by the block till it slides on the plank is ($g=10\ \mathrm{m/s^2}$)
\qn{3} The minimum deviation suffered by a ray emerging through a prism $\left(\mu=\dfrac2{\sqrt3}\right)$ is equal to half the refracting angle of the prism. Find the value of $C$ if the refracting angle of the prism is $\left(\dfrac{2\pi}C\right)$.
\qn{4} It is desired to photograph the image of an object placed at a distance of 3 m from the plane mirror. The camera which is at a distance of 4 m from the mirror should be focussed for a distance of (in m)
\qn{5} The potential energy of a particle of mass 1 kg in motion along the $x$-axis is given by $U=(10-10\cos6x)$ J. If the period of small oscillations is $\dfrac1n$ second, find the value of $n$. $\left(\pi^2\approx10\right)$.
\qn{6} When a plane mirror is placed horizontally on level ground at a distance of 9 m from the foot of a tower, the ray coming from the top of the tower and its reflected ray from the mirror subtends an angle of $90^\circ$. What is the height of the tower ?
\partsec{2}{CHEMISTRY}
\sect{SECTION-I}
\qn{1} Markovnikoff's rule is applicable on which of the following alkene :
\begin{qoptions}
\opt{A}{\qfigi{p2_chem_s1_q01_A}}
\opt{B}{\qfigi{p2_chem_s1_q01_B}}
\opt{C}{\qfigi{p2_chem_s1_q01_C}}
\opt{D}{\qfigi{p2_chem_s1_q01_D}}
\end{qoptions}
\qn{2} Identify reactions correctly matched with their major products ?
\begin{qoptions}
\opt{A}{\qfigi{p2_chem_s1_q02_A}}
\opt{B}{Me--C$\equiv$CH \qfigi{p2_chem_s1_q02_2}}
\opt{C}{Me--C$\equiv$CH \qfigi{p2_chem_s1_q02_3}}
\opt{D}{\qfigi{p2_chem_s1_q02_D}}
\end{qoptions}
\qn{3} Select correct statement(s)
\begin{qoptions}
\opt{A}{Number of second nearest neighbours in primitive cubic unit cell is 12.}
\opt{B}{Minimum inter-ionic distance in $\mathrm{Na_2O}$ is $\sqrt3\,a/2$.}
\opt{C}{Distance between two consecutive hexagonal layer in FCC in $4r\sqrt{2/3}$.}
\opt{D}{Number of next nearest neighbours for $\mathrm{S^{2-}}$ in ZnS type structure is 12.}
\end{qoptions}
\qn{4} At $40^\circ$C, vapour pressure in Torr of methanol and ethanol solution is $P=119x+135$ where $x$ is the mole fraction of methanol. Hence
\begin{qoptions}
\opt{A}{vapour pressure of pure methanol is 254 Torr}
\opt{B}{vapour pressure of pure ethanol is 135 Torr}
\opt{C}{vapour pressure of equimolar mixture of each is 194.5 Torr}
\opt{D}{vapour pressure of pure methanol is 119 Torr}
\end{qoptions}
\qn{5} Lone pair donating power of $\mathrm{(H_3C)N(SiH_3)_2}$ molecule is enhances by :-
\begin{qoptions}
\opt{A}{Replacing only--$\mathrm{CH_3}$ by --$\mathrm{SiH_3}$}
\opt{B}{Replacing one--$\mathrm{SiH_3}$ by --$\mathrm{CH_3}$}
\opt{C}{Replacing both--$\mathrm{SiH_3}$ by --$\mathrm{CH_3}$}
\opt{D}{Replacing both--$\mathrm{SiH_3}$ by --Hydrogen}
\end{qoptions}
\qfullbegin
\qn{6} Consider the following reaction and choose correct statements
\qfig{p2_chem_s1_q06_1}
\qfullsplit
\begin{qoptions}
\opt{A}{Product (B) contain aromatic character.}
\opt{B}{Boiling point of product (C) is less than Helium.}
\opt{C}{All Boron in product (A) are $sp^3$ hybridised.}
\opt{D}{In product (B), $2p_\pi$--$3p\pi$ type back bond present.}
\end{qoptions}
\qfullend
\sect{SECTION-II}
\qn{1} Total number of following carbocation which will undergo rearragement ?
\qfig{p2_chem_s2_q01_1}
\qn{2} \qfigi{p2_chem_s2_q02_1}

If total number of possible stereoisomers for above molecule are $2^n$. The value of $n$ will be :
\qn{3} The normal boiling point of pure water is 373.15 K. Calculate the boiling point (in K) of an aqueous solution containing 11.7 gm NaCl and 34.2 gm sugar ($\mathrm{C_{12}H_{22}O_{11}}$) in 260 gm water. For water, $K_b=0.52$ K--Kg/mol. [Atomic mass of Na = 23, Cl = 35.5]
\qn{4} Henry's law constant for $\mathrm{CO_2}$ in water is $1.6\times10^8$ Pa at 298K. The mass of $\mathrm{CO_2}$ (in gm) is 500g of soda water when packed under 3.2 bar pressure at 298K.
\qn{5} In which of the following molecule $p\pi$ -- $d\pi$ back bonding present?

$\mathrm{CHCl_3}$, $\mathrm{B_3N_3H_6}$, $\mathrm{(CH_3)_3N}$, $\mathrm{SiH_3NCO}$, $\mathrm{:CCl_2}$, $\mathrm{(CH_3)_3SiOH}$, $\mathrm{(SiH_3)_2O}$, $\mathrm{:CF_2}$
\qn{6} Write sum of basicity of oxy acid formed by hydrolysis of :-
\qfig{p2_chem_s2_q06_1}
\qfig{p2_chem_s2_q06_2}
\qfig{p2_chem_s2_q06_3}
\sect{SECTION-III}
\qn{1} How many product(s) is/are formed when trans bute-2-ene reacts with $\mathrm{Br_2}$ / $\mathrm{CCl_4}$
\qn{2} \qfigi{p2_chem_s3_q02_1}

Total number of 1,2 shifts (rearrangements) involve in the above reaction.
\qn{3} In an ionic solid $r_{(+)}=1.6$ \AA\ and $r_{(-)}=1.864$ \AA. Use the radius ratio rule to determine the edge length of the cubic unit cell in \AA. [Given : $\sqrt3=1.732$]
\qn{4} 50 ml of water sample containing temporary hardness only, required 0.1 ml of M/50 HCl for complete neutralisation. Calculate the temporary hardness of water in ppm.
\qn{5} Among the following total number of planar and polar molecules is / are

$\mathrm{Cl_2O,\ P(CH_3)_3,\ N(CH_3)_3,\ I_2Cl_6,\ {}^{\cdot}CH_3,\ NCl_3}$
\qn{6} Find the total number of compound(s), which is/are repelled by magnetic field

$\mathrm{K_2O_2,\ KO_2,\ O_2,\ NO[BF_4],\ KCN,\ Na_2C_2}$
\partsec{3}{MATHEMATICS}
\sect{SECTION-I}
\qn{1} Let $f(x)=\dfrac{2x+1}{x},\ x\neq0$, where \[g(x)=f(f(x))-\dfrac1{f(x)},\quad x\neq0,-\dfrac12,\] then identify the correct statement(s) -
\begin{qoptions}
\opt{A}{$g(x)$ is an odd function}
\opt{B}{$g(x)$ is an even function}
\opt{C}{$g(x)$ is neither even nor odd}
\opt{D}{$g(x)$ is a bounded function}
\end{qoptions}
\qn{2} The values of `$x$' satisfying the equation $\sin^{-1}(x^2-5x+7)=2\tan^{-1}1$ is (are)
\begin{qoptions}
\opt{A}{2}
\opt{B}{3}
\opt{C}{4}
\opt{D}{5}
\end{qoptions}
\qn{3} Consider the equation, $\tan^{-1}(4\{x\})=\tan^{-1}(x+[x])$ where $[\ ]$ and $\{\ \}$ denotes greatest integer function and fractional part function respectively .
\begin{qoptions}
\opt{A}{Number of solutions of the equation are 2.}
\opt{B}{Sum of all real solution(s) of the equation is integer}
\opt{C}{The equation has no real solution}
\opt{D}{The product of all real solution(s) of the equation is non zero.}
\end{qoptions}
\qn{4} Let $a,b,c$ are non-zero real numbers such that $\dfrac1a,\dfrac1b,\dfrac1c$ are in arithmetic progression and $a,b,-2c$ are in geometric progression, then which of the following statement(s) is (are) must be true ?
\begin{qoptions}
\opt{A}{$a^2,b^2,4c^2$ are in geometric progression}
\opt{B}{$-2a,b,-2c$ are in arithmetic progression}
\opt{C}{$a^3+b^3+c^3-3abc=0$}
\opt{D}{$a^2,b^2,c^2$ are in harmonic progression}
\end{qoptions}
\qn{5} Consider a function $f(x)=\dfrac{|x+1|}{x+1}+\dfrac{x-1}{|x-1|}$, then choose correct option(s) regarding function $f$
\begin{qoptions}
\opt{A}{many one}
\opt{B}{odd function}
\opt{C}{bounded function}
\opt{D}{range contain exactly three integers}
\end{qoptions}
\qn{6} Consider the series $S_n=100+95+90+\ldots+T_n$, then correct statement/s is/are (where $T_n$ represents $n^{\mathrm{th}}$ term) -
\begin{qoptions}
\opt{A}{$T_{16}+T_{26}=0$}
\opt{B}{$T_{11}+T_{31}=0$}
\opt{C}{$S_{30}-S_{11}=0$}
\opt{D}{$S_{15}-S_{27}=0$}
\end{qoptions}
\sect{SECTION-II}
\qn{1} Let $f:\mathbb R-\left\{\dfrac12\right\}\to\mathbb R-\left\{\dfrac12\right\}$, $f(x)=\dfrac{x-2}{2x-1}$ be a function such that $x=m$ is the solution of \[f(x)+2f^{-1}(x)+2=f(f(x)),\] then $\dfrac m4$ is equal to
\qn{2} The number of integers in the domain of the function $f(x)=e^{\ell n(\sin^{-1}x)}+e^{\ell n(\cos^{-1}x)}$ is
\qn{3} If the least value of $f(x)=(\tan^{-1}x)^2+(\cot^{-1}x)^2$ where $a$ and $b$ are coprime is $\dfrac{a\pi^2}b$, then find the value of $|a-b|$ ?
\qn{4} If \[\begin{aligned}&2\sin^{-1}\left(\dfrac{9-4x^2}{9+4x^2}\right)+3\cos^{-1}\left(\dfrac{12x}{9+4x^2}\right)\\&\quad+\tan^{-1}\left(\dfrac{12x}{9-4x^2}\right)=\lambda\dfrac\pi2\end{aligned}\] for $x=\dfrac{-17}{13}$, then $\dfrac\lambda2$ is
\qn{5} Sum of first three terms of G. P. with non-integer common ratio is 21 and sum of their squares is 189. If $S$ denotes sum of infinite term of this series, then $\dfrac S4$ is equal to
\qn{6} Minimum possible integral value of the expression $\dfrac{2a^2\sin^4x+(a^4+1)\sin^2x+3a^2}{a^2\sin^2x}$ (wherever defined) is
\columnbreak
\sect{SECTION-III}
\qn{1} Let $f(x)$ be an even function satisfying $f(x)+f(x+2)=f(x+1)\ \forall\ x\in\mathbb R$, where $f(-13)=4$, then $f(-5)+f(1)$ is equal to
\qn{2} If $f(x)=\dfrac{x+a}{x+b}$ and $f:\mathbb R-\{-b\}\to\mathbb R-\{1\}$ is such that $f(x)=f^{-1}(x)$, then value of $|b|$ must be
\qn{3} If $\alpha,\beta$ ($\alpha>\beta$) are two solutions of equation $\tan^{-1}x+\cot^{-1}(-|x|)=2\tan^{-1}6x$, then $2\alpha+3\beta$ is equal to
\qn{4} If range of the function $f(x)=\tan^{-1}(3x^2+(k-1)x+3)$ is $\left[0,\dfrac\pi2\right)$ for $x\in\mathbb R$, then find the sum of all possible real value(s) of $k$ ?
\qn{5} A sequence start with 7120 as its first term and each subsequent term is the sum of cubes of the digits in the previous number. If $N$ is $102^{\mathrm{th}}$ term of this sequence, then sum of digits of $N$ is
\qn{6} Let $a_1=2,\ a_{n+1}=a_n^2-a_n+1$ for $n\geq1$. Find $\displaystyle\sum_{n=1}^{\infty}\dfrac1{a_n}$
\end{multicols}
\end{document}
